\documentclass[11pt,a4paper]{article}

\usepackage[utf8]{inputenc}
\usepackage[T1]{fontenc}
\usepackage[brazilian]{babel}
\usepackage{lmodern}
\usepackage[a4paper,margin=2.5cm]{geometry}
\usepackage{microtype}
\usepackage{hyperref}
\usepackage{enumitem}
\usepackage{booktabs}
\usepackage{array}
\usepackage{tabularx}
\usepackage{graphicx}
\usepackage{amsmath}
\usepackage{tikz}
\usepackage[round]{natbib}
\usepackage{xcolor}

\usetikzlibrary{arrows.meta,positioning}

\definecolor{newblue}{HTML}{1565C0}

\newcommand{\new}[1]{\textcolor{newblue}{#1}}
\newenvironment{addedblock}{\par\color{newblue}}{\par\color{black}}

\hypersetup{
  colorlinks=true,
  linkcolor=black,
  citecolor=black,
  urlcolor=blue
}

\setlength{\parindent}{1.2em}
\setlength{\parskip}{0.4em}
\setlist[itemize]{leftmargin=1.8em}
\setlength{\parindent}{0pt}

\newcolumntype{Y}{>{\raggedright\arraybackslash}X}
\newcolumntype{L}[1]{>{\raggedright\arraybackslash}p{#1}}

\tikzset{
  flowbox/.style={
    draw,
    rounded corners=2pt,
    align=center,
    minimum height=0.95cm,
    minimum width=2.35cm,
    inner sep=4pt
  },
  flowarrow/.style={-{Latex[length=2.5mm]}, thick}
}

\begin{document}

\begin{center}
{\Large Relatório do Projeto Prático}\\[0.4em]
{\large SCC0633 Processamento de Linguagem Natural}\\[0.2em]
{\large SCC5908 Introdução ao Processamento de Língua Natural}\\[0.2em]
\end{center}

\noindent\textit{\small Versão com destaques: trechos em \new{azul} foram adicionados ou reescritos após a revisão da professora.}

\section{Nome do projeto}

\textbf{PLN Core: analisador simbólico de sentimentos para português brasileiro}

\section{Equipe}

\textbf{Nome da equipe:} GLAVJ

\begin{itemize}
  \item João Victor M. Correa (17416061)
  \item Ana Paula de Abreu Batista (12688424)
  \item Guilherme Vaz de Oliveira Taratá (10817476)
  \item Victoria Favero Nunes (15698302)
  \item Leonardo Ferreira Runho (11832958)
\end{itemize}

\begin{addedblock}
\subsection{Autoria do grupo e recursos externos}

A Tabela~\ref{tab:autoria} separa o que foi implementado pela equipe do que foi reaproveitado de recursos externos. Todo o pipeline simbólico, a normalização textual, o tokenizador regex, as regras, o módulo de recomendação, o catálogo de músicas, os testes automatizados e o \emph{harness} de avaliação são autoria do grupo. Para a versão reestruturada da Etapa 1, usada como ponto de comparação da Etapa 2, o único recurso lexical externo ativo é o OpLexicon v3.0.

\begin{table}[h]
\small
\centering
\begin{tabularx}{\linewidth}{@{}L{5.4cm}Y@{}}
\toprule
\textbf{Feito pelo grupo} & \textbf{Recurso externo} \\
\midrule
Pipeline simbólico (\texttt{pipeline.py}) e \emph{factory} de produção (\texttt{factory.py}) & Léxico OpLexicon v3.0 \citep{souza2011construction} \\
Normalização textual e tokenizador regex & --- \\
Regras de negação, intensificação, atenuação, contraste e ênfase, com pesos calibrados manualmente & --- \\
Limiarização e explicação por termos casados & --- \\
Módulo de recomendação (\texttt{recommender.py}) e catálogo curado de 10 músicas (\texttt{recommendations.json}) & --- \\
Testes unitários (\texttt{tests/}) e \emph{harness} de avaliação \emph{Hydra} em \texttt{etapas/etapa1\_simbolica/} & --- \\
\bottomrule
\end{tabularx}
\caption{Separação entre autoria do grupo e recursos externos reaproveitados.}
\label{tab:autoria}
\end{table}
\end{addedblock}

\section{Introdução}

Este projeto tem como objetivo classificar textos curtos em português brasileiro como positivos, negativos ou neutros. A entrada consiste em frases simples, como comentários de redes sociais, e a saída é um rótulo acompanhado de um escore numérico.

A análise de sentimentos é uma tarefa amplamente utilizada em aplicações como monitoramento de opinião pública, avaliação de produtos e apoio à tomada de decisão. Em ambientes digitais, há um grande volume de opiniões textuais produzidas continuamente, o que torna essa tarefa relevante tanto do ponto de vista prático quanto acadêmico.

No contexto do português brasileiro, essa tarefa apresenta desafios adicionais, como variações linguísticas, uso de gírias, ambiguidade semântica e forte dependência de contexto. Entre os principais desafios estão:

\begin{itemize}
  \item \textbf{Dependência de contexto} (ironia e sarcasmo): em ``adorei ficar duas horas esperando no aeroporto'', todas as palavras polarizadas são positivas, mas a intenção é negativa.
  \item \textbf{Ambiguidade lexical}: uma mesma palavra carrega polaridades distintas conforme o contexto \citep{navigli2009wsd}. Em português brasileiro, ``duro'' pode ser físico (neutro), financeiro (``estou duro'', negativo) ou de personalidade (negativo); ``bomba'' pode ser elogio em academia (``bombou no treino'') ou crítica em outro contexto.
  \item \textbf{Negações e intensificadores}: alteram ou amplificam a polaridade dentro de uma janela curta de tokens, como em ``não foi bom'' (inversão) ou ``muito bom'' (ampliação).
  \item \textbf{Cobertura limitada de léxicos}: gírias frequentes em redes sociais (``top'', ``mara'', ``show de bola'', ``irado'') raramente aparecem em léxicos tradicionais, mesmo nos mais amplos como o OpLexicon v3.0 (31\,944 termos).
\end{itemize}

Além disso, trata-se de um problema didaticamente relevante em Processamento de Língua Natural, pois envolve a modelagem de fenômenos linguísticos complexos de forma estruturada. Neste trabalho, adotamos uma abordagem simbólica baseada em léxico e regras explícitas, priorizando interpretabilidade em vez de aprendizado automático \citep{taboada2011lexicon,hutto2014vader,souza2012sentiment}.

\section{Proposta da solução}

\subsection{Arquitetura do pipeline}

A aplicação segue um pipeline simbólico determinístico. Diferente de modelos de aprendizado de máquina, não há treinamento: o processamento ocorre por etapas fixas e interpretáveis.

\begin{figure}[h]
\centering
\resizebox{\linewidth}{!}{
\begin{tikzpicture}[node distance=0.45cm]
  \node[flowbox] (input) {texto\\de entrada};
  \node[flowbox, right=of input] (norm) {normalização};
  \node[flowbox, right=of norm] (tok) {tokenização para texto informal};
  \node[flowbox, right=of tok] (match) {busca no\\léxico};
  \node[flowbox, right=of match] (rules) {regras\\simbólicas};
  \node[flowbox, right=of rules] (sum) {soma dos\\escores};
  \node[flowbox, right=of sum] (label) {rótulo\\final};

  \draw[flowarrow] (input) -- (norm);
  \draw[flowarrow] (norm) -- (tok);
  \draw[flowarrow] (tok) -- (match);
  \draw[flowarrow] (match) -- (rules);
  \draw[flowarrow] (rules) -- (sum);
  \draw[flowarrow] (sum) -- (label);
\end{tikzpicture}
}
\caption{Visão geral da arquitetura da solução.}
\end{figure}

\begin{figure}[h]
\centering
\resizebox{\linewidth}{!}{
\begin{tikzpicture}[node distance=0.42cm]
  \node[flowbox] (base) {termo com\\polaridade};
  \node[flowbox, right=of base] (neg) {negação\\janela de 3 tokens};
  \node[flowbox, right=of neg] (deg) {intensificador\\ou atenuador};
  \node[flowbox, right=of deg] (contrast) {contraste\\antes ou depois de ``mas''};
  \node[flowbox, right=of contrast] (exc) {ênfase por\\exclamação};
  \node[flowbox, right=of exc] (finalscore) {contribuição\\ajustada};

  \draw[flowarrow] (base) -- (neg);
  \draw[flowarrow] (neg) -- (deg);
  \draw[flowarrow] (deg) -- (contrast);
  \draw[flowarrow] (contrast) -- (exc);
  \draw[flowarrow] (exc) -- (finalscore);
\end{tikzpicture}
}
\caption{Fluxo local de ajuste para cada termo encontrado no léxico.}
\end{figure}

\subsection{Recursos e ferramentas}

\begin{addedblock}
A Tabela~\ref{tab:recursos} detalha os recursos ativos no sistema simbólico usado como baseline da Etapa 1: a função, o detalhamento técnico e a origem (própria do grupo ou externa). Implementações alternativas permanecem no repositório para consulta histórica, mas a avaliação oficial desta etapa usa apenas \texttt{oplexicon\_regex}.

\begin{table}[h]
\small
\centering
\begin{tabularx}{\linewidth}{@{}L{2.8cm}L{2.5cm}YL{1.9cm}@{}}
\toprule
\textbf{Recurso} & \textbf{Função} & \textbf{Detalhamento} & \textbf{Origem} \\
\midrule
Normalização textual & Limpeza e padronização do texto bruto & Remove URLs e menções, quebra \emph{hashtags} em \emph{CamelCase} (\texttt{\#MuitoBom} $\rightarrow$ \texttt{muito bom}), comprime caracteres repetidos (\texttt{amooo} $\rightarrow$ \texttt{amo}) e faz \emph{folding} para minúsculas e remoção de acentos. & Grupo \\
Tokenizador regex & Segmentação em \emph{tokens} & Tokenizador determinístico que separa palavras, pontuação e símbolos por expressão regular, sem lematização e sem dependência de modelos externos. & Grupo \\
OpLexicon v3.0 & Polaridade lexical (núcleo) & 31\,944 termos distintos após agregar duplicatas pela média. Cada linha traz \texttt{termo,tipo,polaridade,revisão}; polaridade discreta em $\{-1,0,+1\}$. Baixado e cacheado automaticamente. & PUCRS \citep{souza2011construction} \\
\bottomrule
\end{tabularx}
\caption{Recursos utilizados pelo sistema, com função, detalhamento e origem.}
\label{tab:recursos}
\end{table}
\end{addedblock}

\subsection{Descrição dos processos}

\begin{addedblock}
\textbf{Normalização.}
A camada de normalização para texto informal remove URLs e menções, quebra \emph{hashtags} em \emph{CamelCase}, expande sequências de risada (``kkkkk'', ``rsrs''), comprime caracteres repetidos (``ameiiii'' $\rightarrow$ ``amei'') e preserva \emph{emojis} como \emph{tokens} próprios. Em seguida o texto passa por \emph{folding} (minúsculas e remoção de acentos) para casar com o vocabulário dos léxicos.

\textbf{Tokenização.}
O texto normalizado é segmentado por um tokenizador determinístico baseado em expressões regulares, que separa palavras e pontuação sem lematizar. A renúncia à lematização é deliberada na versão \texttt{oplexicon\_regex}: ela evita dependências morfológicas externas e torna a comparação com os modelos da Etapa 2 mais direta, pois todas as previsões simbólicas decorrem de casamento por forma de superfície no OpLexicon.

\textbf{Léxico ativo.}
O sistema opera com uma única fonte lexical: o OpLexicon v3.0. Entradas duplicadas são agregadas pela média da polaridade, e os termos são normalizados com o mesmo \emph{folding} aplicado ao texto de entrada. Essa decisão deixa explícito o limite do paradigma simbólico usado na comparação: não há treinamento, expansão neural, SentiLex, léxico de gírias ou variantes adicionais no resultado oficial da Etapa 1.

\textbf{Busca no léxico.}
Cada \emph{token} é consultado no OpLexicon normalizado. Quando há correspondência, ele recebe um \emph{escore base} de polaridade, isto é, o valor associado àquele termo antes de qualquer ajuste por regra. Marcadores funcionais (negações e contrastes) são identificados nesta etapa mas não recebem polaridade própria: eles atuam apenas como gatilhos das regras descritas a seguir.

\textbf{Regras simbólicas.}
A Tabela~\ref{tab:regras} resume cada regra com gatilho, janela de aplicação, efeito numérico e exemplo. O arcabouço conceitual remonta a \citet{polanyi2006valenceshifters}, que introduzem o termo \emph{contextual valence shifters} para descrever exatamente esses quatro fenômenos (negação, intensificação, atenuação e contraste). A formulação multiplicativa segue \citet{taboada2011lexicon}, cujo cálculo SO-CAL aplica modificadores percentuais sobre o escore base do termo (\textit{e.g.}, \texttt{very} = $+25\%$, \texttt{deeply} = $+50\%$, \texttt{somewhat} = $-30\%$), e \citet{hutto2014vader}, que padroniza o uso de \emph{booster words} para análise em mídias sociais. A janela de três tokens da negação e a calibração dos modais seguem o estudo empírico de \citet{kiritchenko2016effect}. Os pesos numéricos específicos para o português (\texttt{muito} $1{,}6$, \texttt{super} $1{,}5$, \texttt{pouco} $0{,}6$, contraste $0{,}7/1{,}3$, exclamação $1+0{,}05n$) foram calibrados manualmente pela equipe sobre o \emph{smoke test} de 20 sentenças, situando-se na mesma faixa dos shifters de SO-CAL.

\begin{table}[h]
\small
\centering
\begin{tabularx}{\linewidth}{@{}L{1.9cm}L{3.1cm}L{1.9cm}L{2.6cm}Y@{}}
\toprule
\textbf{Regra} & \textbf{Gatilho} & \textbf{Janela} & \textbf{Efeito} & \textbf{Exemplo} \\
\midrule
Negação & \texttt{não, nem, nunca, jamais, sem} & 3 tokens anteriores & Multiplica por $-1$ (inverte sinal) & ``não foi bom'' $\rightarrow$ \texttt{bom}$\times(-1)$ \\
Intensificação & \texttt{muito} 1{,}6; \texttt{super} 1{,}5; \texttt{demais} 1{,}4; \texttt{bastante} 1{,}3; \texttt{bem}/\texttt{realmente} 1{,}2 & token imediatamente anterior & Multiplica pelo fator do gatilho & ``muito bom'' $\rightarrow$ \texttt{bom}$\times 1{,}6$ \\
Atenuação & \texttt{pouco} 0{,}6; \texttt{meio} 0{,}75; \texttt{quase} 0{,}8 & token imediatamente anterior & Multiplica pelo fator do gatilho & ``pouco útil'' $\rightarrow$ \texttt{útil}$\times 0{,}6$ \\
Contraste & \texttt{mas, porém, contudo, entretanto} (último marcador da frase) & toda a frase & Termos antes: $\times 0{,}7$; depois: $\times 1{,}3$ & ``foi bom, mas péssimo'' \\
Exclamação & ``!'' no texto bruto (até 3) & toda a frase & Multiplica por $1+0{,}05\cdot n$ & ``bom!!'' $\rightarrow$ \texttt{bom}$\times 1{,}10$ \\
\bottomrule
\end{tabularx}
\caption{Regras simbólicas, com gatilho, janela, efeito numérico e exemplo.}
\label{tab:regras}
\end{table}

\textbf{Agregação e classificação.}
O escore final da frase é a soma dos escores ajustados de todos os termos casados. A classificação aplica dois limiares simétricos sobre essa soma:

\begin{equation}
\text{score} = \sum_{i} \text{adjusted}(t_i), \qquad
\text{rótulo} =
\begin{cases}
\text{positivo} & \text{se } \text{score} \geq +0{,}5 \\
\text{negativo} & \text{se } \text{score} \leq -0{,}5 \\
\text{neutro}   & \text{caso contrário}
\end{cases}
\end{equation}

Os limiares simétricos de $\pm 0{,}5$ foram calibrados para textos curtos de mídias sociais, onde o número de termos casados por frase é tipicamente baixo (1 a 3). Limiares mais altos colapsariam a maioria das frases em \emph{neutro}. Quando nenhum termo é encontrado no léxico, o rótulo é \texttt{neutro} por definição.
\end{addedblock}

\subsection{Exemplo de execução}

\begin{addedblock}
A Tabela~\ref{tab:exemplo} ilustra o pipeline completo para a entrada \texttt{O app é muito bom, mas a tela confusa!}. Cada contribuição ajustada é o produto do escore base pelos fatores das regras aplicadas, conforme a Tabela~\ref{tab:regras}.

\begin{table}[h]
\small
\centering
\begin{tabularx}{\linewidth}{@{}L{4cm}Y@{}}
\toprule
\textbf{Etapa} & \textbf{Resultado} \\
\midrule
Entrada & \texttt{O app é muito bom, mas a tela confusa!} \\
Normalização & \emph{folding} para minúsculas, sem acentos; preservação da pontuação e da estrutura da frase \\
Tokenização & \texttt{o, app, e, muito, bom, mas, a, tela, confusa, !} (regex para texto informal, sem lematização) \\
Busca no OpLexicon & \texttt{bom} $(+1{,}0)$ e \texttt{confusa} $(-1{,}0)$ casados no OpLexicon; \texttt{mas} é tratado como marcador funcional de contraste, sem polaridade própria \\
Regras aplicadas & \texttt{bom}: intensificação (\texttt{muito}), pré-contraste, exclamação; \texttt{confusa}: pós-contraste, exclamação \\
Contribuições ajustadas & \texttt{bom}: $1{,}0 \times 1{,}6 \times 0{,}7 \times 1{,}05 \approx +1{,}176$; \texttt{confusa}: $-1{,}0 \times 1{,}3 \times 1{,}05 = -1{,}365$ \\
Agregação e classificação & soma final $\approx -0{,}189$; pelo limiar padrão de $\pm0{,}75$, classificada como \textbf{neutro} \\
\bottomrule
\end{tabularx}
\caption{Aplicação completa do pipeline para \texttt{O app é muito bom, mas a tela confusa!}.}
\label{tab:exemplo}
\end{table}

O exemplo mostra como as regras operam de forma assimétrica em torno de \texttt{mas}: o elogio inicial é amplificado pela intensificação (\texttt{muito}) mas atenuado pelo pré-contraste, enquanto a crítica final é amplificada pelo pós-contraste, dominante para o leitor humano. O sistema captura essa assimetria e devolve \texttt{negativo}, alinhado à interpretação típica de que, após \texttt{mas}, vem o que o falante realmente quis dizer.
\end{addedblock}

\section{Execução}

O repositório remoto do projeto é:

\url{https://github.com/vicmcorrea/PLN-core}

Para executar o projeto, é necessário preparar o ambiente virtual e instalar as dependências.

\subsection*{Execução via Streamlit}

Linux/MacOS:
\begin{verbatim}
uv sync
uv run streamlit run streamlit_app.py
\end{verbatim}

Windows:
\begin{verbatim}
uv sync
uv run streamlit run streamlit_app.py
\end{verbatim}

Para reproduzir a avaliação oficial da Etapa 1 no corpus comum:

\begin{verbatim}
uv run python etapas/etapa1_simbolica/pipelines/run_symbolic_benchmark_suite.py
\end{verbatim}

\subsection*{Recursos adicionais}

Na primeira execução, o OpLexicon v3.0 é automaticamente baixado e armazenado localmente.

Além do \emph{benchmark} via Hydra, o projeto também oferece uma interface web simples utilizando Streamlit:

\begin{verbatim}
streamlit run streamlit_app.py
\end{verbatim}

ou, via opção acessível diretamente no navegador:

\url{https://pln-core.streamlit.app/}


A interface permite visualizar:
\begin{itemize}
  \item rótulo final
  \item escore
  \item texto normalizado
  \item \emph{tokens} gerados para busca no OpLexicon
  \item termos encontrados (com escore base e escore ajustado)
  \item regras simbólicas aplicadas a cada termo
\end{itemize}

\subsection*{Testes}

Para executar os testes atuais do projeto, após o ambiente estar preparado:

\begin{verbatim}
uv run pytest
\end{verbatim}

\begin{addedblock}
\section{Recomendação de músicas}

Como extensão do sistema, foi implementado um módulo de recomendação de músicas baseado no sentimento detectado no texto.

\textbf{Catálogo.} O catálogo, descrito em \texttt{data/recommendations.json}, é um pequeno conjunto curado pela equipe com 10 músicas brasileiras de domínio reconhecido, distribuídas em 4 positivas, 3 negativas e 3 neutras. Cada entrada traz título, artista, identificador no YouTube, rótulo de sentimento (os mesmos três do classificador) e valor de valência. Os rótulos e as valências foram atribuídos manualmente pelo grupo, considerando letra e clima geral da gravação. Nenhum modelo externo, API ou serviço de \emph{streaming} foi utilizado.

\textbf{Valência.} O conceito vem da psicologia do afeto \citep{russell1980circumplex} e é também o nome do atributo de mesma natureza na Spotify Web API. Trata-se de uma medida do grau de positividade emocional de uma faixa: valência alta indica clima alegre e brilhante, valência baixa indica clima triste ou melancólico. Para casar com o escore do classificador, usamos o intervalo $[-1, +1]$. Trata-se de uma aproximação deliberadamente simples, e por isso a rotulagem manual é registrada como limitação na Seção~\ref{sec:conclusao}.

\textbf{Regra de recomendação.} Dado um texto, o sistema filtra músicas com o mesmo rótulo e as ordena pela menor distância entre o escore do texto e a valência da música:

\begin{equation}
d(s) = |\,\text{valence}(s) - \text{score}_{\text{texto}}\,|
\end{equation}

A música recomendada é aquela com menor $d(s)$. O módulo é determinístico e não utiliza aprendizado de máquina.

\textbf{Resultados.} A Tabela~\ref{tab:recomendacoes} mostra a saída real do sistema para quatro entradas representativas, cada uma processada pelo pipeline simbólico ativo (OpLexicon, tokenização regex, regras simbólicas e módulo de recomendação).

\begin{table}[h]
\small
\centering
\begin{tabularx}{\linewidth}{@{}YL{1.6cm}L{1.2cm}L{4.6cm}L{1.4cm}@{}}
\toprule
\textbf{Texto} & \textbf{Rótulo} & \textbf{Escore} & \textbf{Música recomendada} & \textbf{Valência} \\
\midrule
\texttt{Eu amei o filme, foi muito bom!} & positive & $+2{,}06$ & País Tropical (Jorge Ben Jor) & $+0{,}90$ \\
\texttt{Não gostei do app, está bem confuso e bugado.} & negative & $-2{,}00$ & Cálice (Chico Buarque \& Milton Nascimento) & $-0{,}90$ \\
\texttt{O arquivo tem quatro páginas e duas tabelas.} & neutral & $0{,}00$ & Sampa (Caetano Veloso) & $0{,}00$ \\
\texttt{Estou triste com o resultado.} & negative & $-0{,}80$ & Construção (Chico Buarque) & $-0{,}85$ \\
\bottomrule
\end{tabularx}
\caption{Saídas reais do módulo de recomendação para quatro entradas representativas usando o pipeline final.}
\label{tab:recomendacoes}
\end{table}
\end{addedblock}

\begin{addedblock}
\section{Resultados}
\label{sec:resultados}

\textbf{Testes unitários.} O projeto mantém testes automatizados em \texttt{tests/}, cobrindo o pipeline simbólico, os léxicos, os tokenizadores, o módulo de recomendação e a \emph{harness} de avaliação. A execução corrente usa \texttt{uv run pytest}.

\textbf{Protocolo de avaliação.} A avaliação principal foi migrada para o corpus comum que será usado também na Etapa 2: \emph{Portuguese Tweets for Sentiment Analysis}, publicado no Kaggle por \citet{augustopKaggleTweets}. O \emph{split} de treino é \texttt{TrainingDatasets/Train3Classes.csv}, com 100.000 \emph{tweets}, e o \emph{split} de teste é \texttt{TestDatasets/Test3classes.csv}, com 4.999 \emph{tweets}. As três classes são praticamente balanceadas nos dois arquivos. Como o sistema simbólico não aprende parâmetros, o treino é preservado para caracterização do corpus e para a Etapa 2; a métrica oficial da Etapa 1 é calculada no teste. A orquestração usa \emph{Hydra} \citep{yadan2019hydra} em \texttt{etapas/etapa1\_simbolica/pipelines/run\_symbolic\_benchmark\_suite.py}, com saída versionada em \texttt{outputs/etapa1\_symbolic/benchmark\_suite/}. O \emph{run} usado neste relatório é \texttt{20260612\_101202\_400841}.

\begin{table}[h]
\small
\centering
\begin{tabular}{@{}lrrrrrr@{}}
\toprule
\textbf{Split} & \textbf{N} & \textbf{Pos.} & \textbf{Neg.} & \textbf{Neu.} & \textbf{Mediana chars} & \textbf{Mediana tokens} \\
\midrule
Treino & 100.000 & 33.334 & 33.333 & 33.333 & 86 & 13 \\
Teste  & 4.999   & 1.667  & 1.666  & 1.666  & 81 & 11 \\
\bottomrule
\end{tabular}
\caption{Caracterização do corpus comum Kaggle usado por Etapa 1 e Etapa 2.}
\label{tab:dataset-kaggle}
\end{table}

\begin{table}[h]
\small
\centering
\resizebox{\linewidth}{!}{%
\begin{tabular}{@{}llccccc@{}}
\toprule
\textbf{Sistema} & \textbf{Paradigma} & \textbf{Acur.} & \textbf{F1 macro} & \textbf{F1 pos.} & \textbf{F1 neg.} & \textbf{F1 neu.} \\
\midrule
\textbf{\texttt{oplexicon\_regex}} & \textbf{simbólico (OpLexicon + regras)} & \textbf{0,598} & \textbf{0,596} & \textbf{0,652} & \textbf{0,641} & \textbf{0,496} \\
\bottomrule
\end{tabular}%
}
\caption{Resultado oficial da Etapa 1 no \emph{split} de teste Kaggle. Este valor passa a ser o ponto de comparação interno para os modelos da Etapa 2.}
\label{tab:benchmark}
\end{table}

\textbf{Matriz de confusão.} A Tabela~\ref{tab:confusao} mostra para onde os erros do pipeline final vão. O sistema identifica bem polaridades explícitas, especialmente positiva e negativa, mas a classe neutra continua sendo a mais difícil: por definição, ela depende da ausência de evidência lexical e é sensível a termos polarizados usados em manchetes ou contextos informativos.

\begin{table}[h]
\small
\centering
\begin{tabular}{@{}lccc@{}}
\toprule
\textbf{Verdadeiro \textbackslash{} Predito} & \textbf{Positiva} & \textbf{Negativa} & \textbf{Neutra} \\
\midrule
Positiva (1.667) & \textbf{1.133} & 163 & 371 \\
Negativa (1.666) & 229 & \textbf{1.061} & 376 \\
Neutra (1.666)   & 449 & 422 & \textbf{795} \\
\bottomrule
\end{tabular}
\caption{Matriz de confusão do \texttt{oplexicon\_regex} sobre o teste Kaggle.}
\label{tab:confusao}
\end{table}

\textbf{Leitura.} O \texttt{oplexicon\_regex} chega a 0,598 de acurácia e 0,596 de F1 macro no teste Kaggle. A pontuação é coerente com o desenho do método: o sistema não usa rótulos de treino e depende apenas de cobertura lexical e regras locais de composição. O desempenho por classe mostra o limite principal a ser comparado na Etapa 2: positiva e negativa ficam acima de 0,64 de F1, enquanto neutra fica em 0,496. Esse resultado será usado como baseline interpretável para os modelos TF-IDF e neurais no mesmo \emph{split}; portanto, a comparação final do projeto será interna, controlada por corpus, \emph{split} e métrica, em vez de depender de resultados externos com protocolos diferentes.
\end{addedblock}

\begin{addedblock}
\section{Conclusão}
\label{sec:conclusao}

O presente trabalho apresentou um sistema simbólico para análise de sentimentos em português brasileiro, baseado no OpLexicon v3.0 e em regras explícitas, priorizando interpretabilidade ao longo de todo o pipeline. A versão oficial reestruturada para a comparação com a Etapa 2 é o \texttt{oplexicon\_regex}: OpLexicon, normalização textual, tokenização regex e regras de negação, intensificação, atenuação, contraste, caixa alta e ênfase exclamativa.

A avaliação sobre o corpus Kaggle comum fixa o que esse paradigma entrega antes de treinamento estatístico: 0,598 de acurácia e 0,596 de F1 macro em 4.999 \emph{tweets} de teste balanceados. As limitações são as esperadas em sistemas simbólicos: dependência da cobertura lexical, simplificação de fenômenos complexos via regras heurísticas e análise predominantemente local do contexto. O módulo de recomendação, embora funcional, baseia-se em valores de valência atribuídos manualmente pela equipe, o que limita a precisão das sugestões fora do catálogo curado.

Os caminhos naturais de evolução já se desenham a partir da análise dos erros: a classe neutra é a que mais sofre (F1 0,496), porque qualquer ruído lexical a empurra para as caudas positiva ou negativa. A Etapa 2 parte exatamente desse ponto: treinar classificadores TF-IDF e modelos neurais no mesmo \emph{split} de treino e compará-los com esta linha simbólica, mantendo corpus, métricas e artefatos de saída controlados.
\end{addedblock}

\bibliographystyle{plainnat}
\bibliography{references}

\end{document}
